% Banque d'exercices — chapitre 5
% Le chapitre est déterminé par le nom de ce fichier.

\begin{exo}[Chauffage isochore puis compression isotherme]{chauffage-isochore-gp-mono}
\seoready{true}
\keywords{gaz parfait monoatomique, chauffage isochore, compression isotherme, énergie interne, travail, chaleur, premier principe}

Une mole de gaz parfait monoatomique subit deux transformations successives :
\begin{enumerate}
\item un chauffage à volume constant de $T_A=300\ \mathrm{K}$ à
$T_B=500\ \mathrm{K}$ ;
\item une compression isotherme, à $T_B$, du volume $V_B$ au volume
$V_C=V_B/2$.
\end{enumerate}

Calculer simplement $W$, $\Delta U$ et $Q$ pour chaque transformation. On
prendra $R=8{,}314\ \mathrm{J\,mol^{-1}\,K^{-1}}$.

\begin{indication}
À volume constant, le travail est nul. Pour la compression isotherme, utiliser
$W=-nRT_B\ln(V_C/V_B)$.
\end{indication}

\begin{solution}
\textbf{1. Chauffage isochore.} Le volume est constant, donc
\[
W_{AB}=0.
\]
Avec $n=1\ \mathrm{mol}$ et $\Delta T=200\ \mathrm{K}$,
\[
\Delta U_{AB}
=\frac32R\Delta T
=\frac32\times8{,}314\times200
\simeq2{,}49\ \mathrm{kJ}.
\]
Le premier principe donne
\[
Q_{AB}=\Delta U_{AB}-W_{AB}\simeq2{,}49\ \mathrm{kJ}.
\]
\textbf{2. Compression isotherme.} La température reste égale à
$T_B=500\ \mathrm{K}$, donc
\[
W_{BC}
=-RT_B\ln\!\left(\frac{V_C}{V_B}\right)
=RT_B\ln2
\simeq2{,}88\ \mathrm{kJ}.
\]
Pour un gaz parfait, l'énergie interne ne dépend que de la température :
\[
\Delta U_{BC}=0.
\]
Ainsi,
\[
Q_{BC}=-W_{BC}\simeq-2{,}88\ \mathrm{kJ}.
\]
Pendant la compression, le gaz reçoit du travail et cède la même énergie sous
forme de chaleur.
\end{solution}
\end{exo}

\begin{exo}[Travail et chaleur entre deux états par deux chemins]{travail-chaleur-deux-chemins}
\keywords{travail des forces de pression, chaleur, isotherme, isochore, diagramme de Clapeyron, différentielle exacte, gaz parfait}

On considère $n$ moles de gaz parfait monoatomique,
$U=\frac{3}{2}nRT$. On joint l'état $A(T_1,V_1)$ à l'état $B(T_2,V_2)$
par deux chemins quasi-statiques, avec $P_{\mathrm{ext}}=P$ à chaque instant :
\begin{itemize}
\item $\Gamma_1$ : une détente isotherme à $T_1$, de $V_1$ à $V_2$,
puis une transformation isochore à $V_2$, de $T_1$ à $T_2$ ;
\item $\Gamma_2$ : une transformation isochore à $V_1$, de $T_1$ à $T_2$,
puis une détente isotherme à $T_2$, de $V_1$ à $V_2$.
\end{itemize}

\begin{enumerate}
\item Dans le cas $T_2>T_1$ et $V_2>V_1$, représenter les chemins
$\Gamma_1$ et $\Gamma_2$ dans le diagramme de Clapeyron $(P,V)$. Indiquer
les états intermédiaires et le sens de parcours.
\item Calculer le travail reçu $W_1$ et $W_2$ le long de chacun des deux
chemins. Que vaut leur différence ?
\item Calculer $\Delta U$ le long de chacun des deux chemins, puis en déduire
$Q_1$ et $Q_2$. Conclure sur la nature de $\delta Q$.
\item Application numérique : $n=1$ mol, $T_1=300$ K, $T_2=600$ K et
$V_2=2V_1$.
\end{enumerate}

\begin{indication}
Dans le diagramme de Clapeyron, une isochore est verticale et une isotherme
de gaz parfait vérifie $P=nRT/V$. L'isotherme de température $T_2$ est située
au-dessus de celle de température $T_1$. Une isochore ne contribue pas au
travail des forces de pression.
\end{indication}

\begin{solution}
\textbf{Question 1.} Notons $C$ l'état intermédiaire du chemin $\Gamma_1$
et $D$ celui du chemin $\Gamma_2$. Leurs pressions et leurs volumes sont
\[
\begin{aligned}
A&=\left(\frac{nRT_1}{V_1},V_1\right),&
C&=\left(\frac{nRT_1}{V_2},V_2\right),\\
D&=\left(\frac{nRT_2}{V_1},V_1\right),&
B&=\left(\frac{nRT_2}{V_2},V_2\right).
\end{aligned}
\]
Puisque $T_2>T_1$, l'isotherme $T_2$ est au-dessus de l'isotherme $T_1$.
Le chemin $\Gamma_1$ se trace de $A$ vers $C$, vers la droite le long de
l'isotherme inférieure $T_1$, puis verticalement de $C$ vers $B$ le long de
l'isochore $V_2$. Le chemin $\Gamma_2$ monte d'abord verticalement de $A$
vers $D$ le long de l'isochore $V_1$, puis va vers la droite de $D$ vers $B$
le long de l'isotherme supérieure $T_2$. Les flèches des deux chemins sont
donc
\[
\Gamma_1:A\xrightarrow{\,T=T_1\,}C\xrightarrow{\,V=V_2\,}B,
\qquad
\Gamma_2:A\xrightarrow{\,V=V_1\,}D\xrightarrow{\,T=T_2\,}B.
\]

\textbf{Question 2.} Avec la convention du travail reçu par le gaz,
$\delta W=-P\,\mathrm{d}V$. Une isochore ne contribue pas et, sur une
isotherme à la température $T$,
\[
W=-\int_{V_1}^{V_2}\frac{nRT}{V}\,\mathrm{d}V
=-nRT\ln\frac{V_2}{V_1}.
\]
On obtient donc
\[
W_1=-nRT_1\ln\frac{V_2}{V_1},
\qquad
W_2=-nRT_2\ln\frac{V_2}{V_1},
\]
puis
\[
W_1-W_2
=nR(T_2-T_1)\ln\frac{V_2}{V_1}\neq 0.
\]
Les états initial et final sont identiques, mais les travaux diffèrent :
$\delta W$ n'est pas une différentielle exacte.

\textbf{Question 3.} L'énergie interne du gaz parfait monoatomique ne dépend
que de la température. Sa variation est donc la même sur les deux chemins :
\[
\Delta U=U_B-U_A=\frac{3}{2}nR(T_2-T_1).
\]
Le premier principe, $\Delta U=W+Q$, donne $Q=\Delta U-W$, d'où
\[
\begin{aligned}
Q_1&=\frac{3}{2}nR(T_2-T_1)
     +nRT_1\ln\frac{V_2}{V_1},\\
Q_2&=\frac{3}{2}nR(T_2-T_1)
     +nRT_2\ln\frac{V_2}{V_1}.
\end{aligned}
\]
Ainsi,
\[
Q_2-Q_1=nR(T_2-T_1)\ln\frac{V_2}{V_1}\neq 0.
\]
La chaleur échangée dépend elle aussi du chemin : $\delta Q$ n'est pas une
différentielle exacte, contrairement à $\mathrm{d}U$.

\textbf{Question 4.} Avec $\ln 2\simeq0{,}693$,
\[
\begin{aligned}
W_1&\simeq-1{,}73\ \mathrm{kJ},&
W_2&\simeq-3{,}46\ \mathrm{kJ},\\
\Delta U&\simeq3{,}74\ \mathrm{kJ},&
Q_1&\simeq5{,}47\ \mathrm{kJ},&
Q_2&\simeq7{,}20\ \mathrm{kJ}.
\end{aligned}
\]
Le gaz fournit deux fois plus de travail au cours de $\Gamma_2$ et reçoit
la chaleur supplémentaire correspondante.
\end{solution}
\end{exo}

\begin{exo}[Cycle triangulaire]{cycle-triangulaire-diagonale}
\keywords{cycle moteur, transformation linéaire, diagramme de Clapeyron, aire du cycle, premier principe}

On considère $n$ moles de gaz parfait, de capacité thermique molaire $c_V$
constante, décrivant le cycle quasi-statique $A\to B\to C\to A$, avec
\[
A=(P_0,V_0),
\qquad
B=(2P_0,V_0),
\qquad
C=(P_0,2V_0).
\]
La branche $A\to B$ est isochore, la branche $B\to C$ est un segment de droite
dans le plan $(P,V)$ et la branche $C\to A$ est isobare.

\begin{enumerate}
\item Représenter le cycle dans le diagramme de Clapeyron et préciser son sens
de parcours. Est-il moteur ou récepteur ?
\item Déterminer l'équation $P(V)$ de la droite reliant $B$ à $C$.
\item En posant $T_0=P_0V_0/(nR)$, déterminer les températures $T_A$, $T_B$ et
$T_C$. Pourquoi $\Delta U_{B\to C}$ est-elle nulle alors que cette branche
n'est pas isotherme ?
\item Calculer le travail reçu sur chacune des trois branches, puis sur le cycle
complet. Retrouver la valeur absolue du travail total à partir de l'aire du
triangle.
\item Calculer $\Delta U$ et $Q$ sur chaque branche, puis vérifier que
$\Delta U_{\mathrm{cycle}}=0$ et $Q_{\mathrm{cycle}}=-W_{\mathrm{cycle}}$.
\end{enumerate}

\begin{indication}
Sur $B\to C$, chercher $P(V)=aV+b$ en imposant le passage par $B$ et $C$.
Pour calculer une variation d'énergie interne, seules les températures aux
extrémités de la branche sont nécessaires.
\end{indication}

\begin{solution}
\textbf{Questions 1--2.} Le cycle est parcouru dans le sens horaire : il est
moteur. L'équation de la droite $B\to C$ est
\[
P(V)=P_0\left(3-\frac{V}{V_0}\right).
\]

\textbf{Question 3.} L'équation d'état donne
\[
T_A=T_0,
\qquad
T_B=T_C=2T_0.
\]
La température varie le long de $B\to C$, mais elle reprend la même valeur à
ses deux extrémités. Comme l'énergie interne d'un gaz parfait ne dépend que de
la température à quantité de matière fixée, $\Delta U_{B\to C}=0$.

\textbf{Question 4.} Les travaux reçus valent
\[
W_{AB}=0,
\qquad
W_{BC}=-\int_{V_0}^{2V_0}P(V)\,dV=-\frac32P_0V_0,
\qquad
W_{CA}=P_0V_0.
\]
Ainsi,
\[
\boxed{W_{\mathrm{cycle}}=-\frac12P_0V_0}.
\]
Sa valeur absolue est bien l'aire du triangle, de base $V_0$ et de hauteur
$P_0$.

\textbf{Question 5.} Puisque $nRT_0=P_0V_0$,
\[
\begin{array}{c|c|c|c}
\text{Branche} & W & \Delta U & Q=\Delta U-W\\ \hline
A\to B & 0 & \dfrac{c_V}{R}P_0V_0
& \dfrac{c_V}{R}P_0V_0\\[4pt]
B\to C & -\dfrac32P_0V_0 & 0 & \dfrac32P_0V_0\\[4pt]
C\to A & P_0V_0 & -\dfrac{c_V}{R}P_0V_0
& -\left(\dfrac{c_V}{R}+1\right)P_0V_0
\end{array}
\]
En sommant les trois branches, on obtient
\[
\boxed{
\Delta U_{\mathrm{cycle}}=0,
\qquad
Q_{\mathrm{cycle}}=\frac12P_0V_0=-W_{\mathrm{cycle}}.
}
\]
\end{solution}
\end{exo}

\begin{exo}[Chauffage à pression constante : découverte de $C_P$]{decouverte-cp}
\keywords{capacité thermique, pression constante, relation de Mayer, chauffage isobare, gaz parfait}

Une mole de gaz parfait, de capacité thermique molaire $C_V$, est enfermée
dans un cylindre de volume initial $V_0$. Le cylindre est fermé par un piston
mobile sans frottement, soumis à une pression extérieure constante $P_0$.
Une résistance électrique apporte au gaz une chaleur $Q>0$ de façon
quasi-statique.

\begin{enumerate}
\item Faire un schéma simple du dispositif. Indiquer le gaz, le piston, la
résistance, la pression extérieure $P_0$ et les échanges d'énergie.
\item Sans utiliser de valeurs numériques, exprimer la température initiale
$T_0$, puis déterminer l'état final $(P_f,V_f,T_f)$, le travail $W$ reçu par
le gaz et sa variation d'énergie interne $\Delta U$. Quelle fraction de $Q$
sert à repousser l'atmosphère ?
\item En déduire la capacité thermique molaire à pression constante
$C_P=(\delta Q/\mathrm{d}T)_P$ en fonction de $C_V$ et de $R$.
\item Application numérique pour un gaz parfait monoatomique :
$C_V=\frac32R$, $V_0=25\ \mathrm{L}$,
$P_0=1{,}0\times10^5\ \mathrm{Pa}$, $Q=1000\ \mathrm{J}$ et
$R=8{,}314\ \mathrm{J\,mol^{-1}\,K^{-1}}$.
\end{enumerate}

\begin{indication}
La transformation est isobare : $W=-P_0\,\Delta V$. Pour une mole de gaz
parfait, $P_0\,\Delta V=R\,\Delta T$. Injecter ces relations dans le premier
principe.
\end{indication}

\begin{solution}
\textbf{Question 1.} Le schéma doit montrer le gaz dans un cylindre fermé par
un piston mobile. Une flèche entrante représente la chaleur $Q$ apportée par
la résistance. La pression extérieure $P_0$ s'exerce sur le piston, qui se
déplace vers le haut lorsque le gaz se dilate.

\textbf{Question 2 --- Calcul formel.} L'équation d'état initiale donne
\[
T_0=\frac{P_0V_0}{R}.
\]
La pression reste constante, donc $P_f=P_0$. De plus,
$W=-P_0\Delta V=-R\Delta T$. Le premier principe donne alors
\[
\Delta U=C_V\Delta T=Q+W=Q-R\Delta T,
\]
d'où
\[
\Delta T=\frac{Q}{C_V+R}.
\]
L'état final est donc
\[
P_f=P_0,\qquad
T_f=T_0+\frac{Q}{C_V+R},\qquad
V_f=V_0+\frac{RQ}{P_0(C_V+R)}.
\]
Enfin,
\[
W=-\frac{RQ}{C_V+R},
\qquad
\Delta U=\frac{C_VQ}{C_V+R}.
\]
La fraction de la chaleur servant à repousser l'atmosphère vaut
\[
\frac{-W}{Q}=\frac{R}{C_V+R}.
\]

\textbf{Question 3.} Comme $Q=(C_V+R)\Delta T$ à pression constante,
\[
C_P=C_V+R.
\]
C'est la relation de Mayer.

\textbf{Question 4 --- Application numérique.} Pour un gaz monoatomique,
\[
C_P=\frac52R\simeq20{,}8\ \mathrm{J\,mol^{-1}\,K^{-1}}.
\]
On obtient
\[
T_0\simeq301\ \mathrm{K},\qquad
\Delta T\simeq48{,}1\ \mathrm{K},\qquad
T_f\simeq349\ \mathrm{K},\qquad
V_f=29{,}0\ \mathrm{L}.
\]
Enfin,
\[
W=-400\ \mathrm{J},\qquad
\Delta U=600\ \mathrm{J}.
\]
La fraction de chaleur utilisée pour repousser l'atmosphère est
$-W/Q=2/5$, soit $40\,\%$.
\end{solution}
\end{exo}

\begin{exo}[Adiabatique réversible d'un gaz parfait : lois de Laplace]{lois-de-laplace}
\keywords{loi de Laplace, adiabatique réversible, indice adiabatique, gamma, détente adiabatique, premier principe différentiel}

On considère $n$ moles de gaz parfait d'indice adiabatique $\gamma = C_P/C_V$ supposé constant. On admet la relation de Mayer $C_P - C_V = nR$ (voir l'exercice sur la découverte de $C_P$).

\begin{enumerate}
\item Montrer que
\[
C_V = \frac{nR}{\gamma - 1}, \qquad C_P = \frac{\gamma\, nR}{\gamma - 1}.
\]
\item Le gaz subit une transformation adiabatique réversible. En partant du premier principe sous forme différentielle, montrer que
\[
\frac{\mathrm{d}T}{T} + (\gamma - 1)\frac{\mathrm{d}V}{V} = 0,
\]
puis en déduire que $TV^{\gamma-1} = \text{cste}$ le long de la transformation.
\item En déduire les deux autres formes de la loi de Laplace : $PV^\gamma = \text{cste}$ et $T^\gamma P^{1-\gamma} = \text{cste}$.
\item Détente adiabatique réversible qui double le volume, pour un gaz monoatomique ($\gamma = 5/3$), $n = 1$ mol, $T_I = 300$ K. Calculer $T_F$ puis $W$, et commenter les signes.
\item Dans le diagramme de Clapeyron $(P,V)$, comparer en un même point la pente d'une adiabatique réversible et celle d'une isotherme.
\end{enumerate}

\begin{indication}
Pour la question 2 : $\mathrm{d}U = C_V\,\mathrm{d}T$ et $\delta W = -P\,\mathrm{d}V$ avec $P = nRT/V$, puis séparer les variables. Pour la question 5, différencier $PV^\gamma = \text{cste}$ et $PV = \text{cste}$.
\end{indication}

\begin{solution}
\textbf{Question 1.} De $C_P = \gamma C_V$ et $C_P - C_V = nR$ : $C_V(\gamma - 1) = nR$, d'où les deux formules.

\textbf{Question 2.} Adiabatique : $\delta Q = 0$, donc $\mathrm{d}U = \delta W$, soit $C_V\,\mathrm{d}T = -P\,\mathrm{d}V = -\dfrac{nRT}{V}\mathrm{d}V$. Avec $C_V = nR/(\gamma-1)$ :
\[
\frac{\mathrm{d}T}{T} = -(\gamma-1)\frac{\mathrm{d}V}{V}.
\]
En intégrant : $\ln T + (\gamma-1)\ln V = \text{cste}$, c'est-à-dire $TV^{\gamma-1} = \text{cste}$.

\textbf{Question 3.} En substituant $T = PV/(nR)$ : $PV^\gamma = \text{cste}$. En substituant $V = nRT/P$ : $T^\gamma P^{1-\gamma} = \text{cste}$.

\textbf{Question 4.} $T_F = T_I\,2^{-(\gamma-1)} = 300 \times 2^{-2/3} \approx 189$ K. Comme $Q = 0$ :
\[
W = \Delta U = \tfrac{3}{2}R(T_F - T_I) \approx \tfrac{3}{2}\times 8{,}314 \times (-111) \approx -1{,}4\ \text{kJ}.
\]
$W < 0$ : le gaz fournit du travail à l'extérieur. Ne pouvant le prélever sur aucune source de chaleur ($Q=0$), il puise dans son énergie interne : d'où le refroidissement $\Delta T < 0$.

\textbf{Question 5.} En différenciant $PV = \text{cste}$ : $\left(\frac{\mathrm{d}P}{\mathrm{d}V}\right)_{\text{isoT}} = -\dfrac{P}{V}$. En différenciant $PV^\gamma = \text{cste}$ : $\left(\frac{\mathrm{d}P}{\mathrm{d}V}\right)_{\text{adiab}} = -\gamma\dfrac{P}{V}$. L'adiabatique est $\gamma$ fois plus pentue que l'isotherme : comprimer sans évacuer la chaleur fait monter la pression plus vite, car la température augmente aussi.
\end{solution}
\end{exo}

\begin{exo}[Travail d'une transformation polytropique $PV^k=\text{cste}$]{transformation-polytropique}
\keywords{transformation polytropique, travail des forces de pression, adiabatique reversible, loi de Laplace}

On considere une transformation quasi-statique d'un gaz parfait telle que $PV^k=\text{cste}$, $k\in\mathbb{R}$, $k\neq1$.

\begin{enumerate}
\item Calculer $W_{A\to B}$ en fonction de $P_A,V_A,P_B,V_B$ et $k$.
\item Retrouver le resultat pour une isobare ($k=0$).
\item Conclure directement pour une isochore sans utiliser la formule generale.
\item Expliquer pourquoi $k=1$ est exclu, puis retrouver $W=-nRT\ln(V_B/V_A)$.
\item Pour $k=\gamma=5/3$ (adiabatique reversible, monoatomique), montrer que $W=\Delta U$.
\end{enumerate}

\begin{indication}
Poser $P(V)=CV^{-k}$ avec $C=P_AV_A^k$, integrer, puis utiliser $CV^{1-k}=PV$.
\end{indication}

\begin{solution}
\textbf{Question 1.}
\[
W = \frac{P_AV_A - P_BV_B}{1-k}
\]

\textbf{Question 2.} $k=0$, $P_A=P_B=P$ : $W=P(V_A-V_B)=-P\Delta V$. \checkmark

\textbf{Question 3.} Isochore : $\mathrm{d}V=0$ a chaque instant, donc $W=0$. La formule polytropique ne s'applique pas ($k\to\infty$).

\textbf{Question 4.} Pour $k=1$ le facteur $1/(1-k)$ diverge (primitive logarithmique). Avec $PV=nRT=\text{cste}$ :
\[
W = -nRT\ln\!\tfrac{V_B}{V_A}.
\]

\textbf{Question 5.} $k=5/3$ :
\[
W = \tfrac{3}{2}(P_BV_B-P_AV_A) = \tfrac{3}{2}nR(T_B-T_A) = \Delta U.
\]
Coherent avec le premier principe ($Q=0$). \checkmark
\end{solution}
\end{exo}


\begin{exo}[Résistance chauffante dans un cylindre cloisonné]{resistance-chauffante-cloison}
\keywords{premier principe, piston adiabatique, loi de Laplace, compression adiabatique réversible, bilan d'énergie, hélium}

Un cylindre horizontal rigide, aux parois calorifugées, est divisé en deux compartiments $C_1$ et $C_2$ par un piston mobile sans frottement, lui aussi calorifugé. Chaque compartiment contient de l'hélium (gaz parfait monoatomique, $\gamma = 5/3$) dans le même état initial : $P_0 = 1{,}0\times10^5$ Pa, $V_0 = 1{,}0$ L, $T_0 = 300$ K. Une résistance électrique placée dans $C_1$ lui apporte de la chaleur très lentement, jusqu'à ce que la pression commune atteigne $P_1 = 1{,}5\times10^5$ Pa. On supposera toutes les transformations quasi-statiques.

\begin{enumerate}
\item Quelle transformation subit le gaz de $C_2$ ? En déduire son volume $V_2$ et sa température $T_2$ finaux.
\item En déduire le volume $V_1$ et la température $T_1$ finaux du gaz de $C_1$.
\item Calculer les variations d'énergie interne $\Delta U_1$ et $\Delta U_2$. On remarquera que pour un gaz parfait monoatomique, $U = \frac{3}{2}PV$.
\item En appliquant le premier principe au système global $\{C_1 + C_2\}$, déterminer la chaleur $Q$ fournie par la résistance. Vérifier la cohérence en appliquant le premier principe à $C_2$ seul, puis à $C_1$ seul.
\end{enumerate}

\begin{indication}
Le gaz de $C_2$ n'échange aucune chaleur (piston et parois calorifugés) et est comprimé de façon quasi-statique : c'est une adiabatique réversible, la loi de Laplace $PV^\gamma = \text{cste}$ s'applique. Le volume total $V_1 + V_2 = 2V_0$ est fixé. Pour le système global, les parois externes sont rigides : seul $Q$ entre.
\end{indication}

\begin{solution}
\textbf{Question 1.} Le gaz de $C_2$ est comprimé lentement, sans frottement et sans échange de chaleur : il subit une compression adiabatique réversible. Loi de Laplace :
\[
V_2 = V_0\!\left(\frac{P_0}{P_1}\right)^{1/\gamma} = 1{,}0\times\left(\frac{1}{1{,}5}\right)^{3/5} \approx 0{,}78\ \text{L},
\qquad
T_2 = T_0\,\frac{P_1V_2}{P_0V_0} \approx 353\ \text{K}.
\]

\textbf{Question 2.} Le cylindre est rigide : $V_1 = 2V_0 - V_2 \approx 1{,}22$ L, et l'équation d'état donne
\[
T_1 = T_0\,\frac{P_1V_1}{P_0V_0} = 300 \times 1{,}5 \times 1{,}22 \approx 547\ \text{K}.
\]

\textbf{Question 3.} Avec $U = \tfrac{3}{2}PV$ :
\[
\Delta U_1 = \tfrac{3}{2}(P_1V_1 - P_0V_0) \approx \tfrac{3}{2}(182{,}4 - 100) \approx 124\ \text{J},
\qquad
\Delta U_2 = \tfrac{3}{2}(P_1V_2 - P_0V_0) \approx \tfrac{3}{2}(117{,}6 - 100) \approx 26\ \text{J}.
\]

\textbf{Question 4.} Le système global n'échange aucun travail (parois externes rigides) et reçoit uniquement $Q$ :
\[
Q = \Delta U_1 + \Delta U_2 = \tfrac{3}{2}(P_1 - P_0)(2V_0) = \tfrac{3}{2}\times 0{,}5\times10^5 \times 2\times10^{-3} = 150\ \text{J}.
\]
Cohérence : $C_2$ est adiabatique, donc $W_2 = \Delta U_2 \approx 26$ J est le travail que lui fournit le piston. Le gaz de $C_1$ reçoit $Q$ de la résistance et cède ce travail au piston : $\Delta U_1 = Q - W_2 \approx 150 - 26 = 124$ J. \checkmark

On note l'efficacité du bilan global : nul besoin de connaître le détail des échanges internes pour trouver $Q$.
\end{solution}
\end{exo}

\begin{exo}[Cycle isotherme--isobare--isochore : bilan de synthèse]{travail-trois-chemins}
\keywords{travail des forces de pression, isotherme, isobare, isochore, diagramme de Clapeyron, premier principe}

On considère un système fermé constitué de $n$ moles de gaz parfait monoatomique ($U = \frac{3}{2}nRT$), initialement dans l'état $A(P_A, V_A, T_A)$. Toutes les transformations sont quasi-statiques.

\begin{enumerate}
\item Compression isotherme $A\to B$ jusqu'à $V_B = V_A/2$. Calculer $P_B$, puis $W$, $\Delta U$ et $Q$ sur cette branche.
\item Détente isobare $B\to C$ jusqu'à $V_C = V_A$. Calculer $T_C$, puis $W$, $\Delta U$ et $Q$.
\item Refroidissement isochore $C\to A$ qui referme le cycle. Calculer $W$, $\Delta U$ et $Q$.
\item Représenter le cycle $A\to B\to C\to A$ dans le diagramme de Clapeyron $(P,V)$. Application numérique pour $n = 1$ mol, $P_A = 1{,}0$ bar, $T_A = 300$ K : donner $(P,V,T)$ en chacun des trois points.
\item Comparer $W_{A\to B}$ au travail $W_{A\to C\to B}$ reçu le long du chemin isochore $A\to C$ puis isobare $C\to B$. Commenter.
\item Le cycle $A\to B\to C\to A$ est-il moteur ou récepteur ? Calculer le travail total et vérifier la cohérence avec le sens de parcours.
\item Représenter le même cycle dans le diagramme $(T,P)$.
\end{enumerate}

\begin{indication}
Sur une isotherme de gaz parfait, remplacer $P$ par $nRT/V$ dans $\delta W = -P\,\mathrm{d}V$ avant d'intégrer. En diagramme $(P,V)$, l'aire du cycle vaut le travail fourni si le parcours est horaire. Convention banquier : tout ce qui entre dans le système est compté positivement.
\end{indication}

\begin{solution}
\textbf{Question 1.} Isotherme : $T_B = T_A$, et $P_AV_A = P_BV_B$ donne $P_B = 2P_A$.
\[
W_{AB} = -\int_{V_A}^{V_A/2}\frac{nRT_A}{V}\,\mathrm{d}V = -nRT_A\ln\frac{V_B}{V_A} = nRT_A\ln 2 > 0.
\]
$U$ ne dépend que de $T$ : $\Delta U = 0$, donc $Q_{AB} = -W_{AB} = -nRT_A\ln 2 < 0$ (le gaz évacue la chaleur de compression).

\textbf{Question 2.} $P_CV_C = 2P_AV_A = 2nRT_A$ donc $T_C = 2T_A$. À pression constante $P_B = 2P_A$ :
\[
W_{BC} = -2P_A\left(V_A - \tfrac{V_A}{2}\right) = -P_AV_A = -nRT_A, \qquad \Delta U_{BC} = \tfrac{3}{2}nR(T_C - T_A) = \tfrac{3}{2}nRT_A,
\]
d'où $Q_{BC} = \Delta U_{BC} - W_{BC} = \tfrac{5}{2}nRT_A$.

\textbf{Question 3.} Isochore : $W_{CA} = 0$ et $Q_{CA} = \Delta U_{CA} = \tfrac{3}{2}nR(T_A - 2T_A) = -\tfrac{3}{2}nRT_A$.

\textbf{Question 4.} $V_A = RT_A/P_A \approx 24{,}9$ L. Donc $A$ : $(1$ bar, $24{,}9$ L, $300$ K$)$, $B$ : $(2$ bar, $12{,}5$ L, $300$ K$)$, $C$ : $(2$ bar, $24{,}9$ L, $600$ K$)$.

\textbf{Question 5.} $W_{A\to C\to B} = 0 + \left[-2P_A\left(\tfrac{V_A}{2} - V_A\right)\right] = P_AV_A = nRT_A \approx 2494$ J, alors que $W_{A\to B} = nRT_A\ln 2 \approx 1729$ J. Mêmes états extrêmes, travaux différents : le travail dépend du chemin suivi, $\delta W$ n'est pas une différentielle exacte.

\textbf{Question 6.} $W_{\text{total}} = nRT_A\ln 2 - nRT_A = nRT_A(\ln 2 - 1) \approx -765$ J $< 0$ : le gaz fournit du travail, le cycle est moteur. C'est cohérent : en diagramme $(P,V)$, le parcours est horaire (isotherme en bas, isobare en haut) et l'aire du cycle vaut le travail fourni.

\textbf{Question 7.} En diagramme $(T,P)$ : $A\to B$ est un segment vertical à $T = T_A$ (de $P_A$ à $2P_A$) ; $B\to C$ un segment horizontal à $P = 2P_A$ (de $T_A$ à $2T_A$) ; $C\to A$ un segment de la droite passant par l'origine $P = (nR/V_A)\,T$ (isochore).
\end{solution}
\end{exo}

\begin{exo}[Cycle rectangulaire biisobare-biisochore : bilan complet]{cycle-biisobare-biisochore}
\seoready{true}
\keywords{cycle thermodynamique, diagramme de Clapeyron, travail, chaleur, aire, gaz parfait, cycle moteur}

On considère $n$ moles de gaz parfait monoatomique, avec $U=\frac{3}{2}nRT$, parcourant quasi-statiquement le cycle rectangulaire horaire
\[
A(P_1,V_1)\to B(P_2,V_1)\to C(P_2,V_2)
\to D(P_1,V_2)\to A,
\]
où $V_2>V_1$ et $P_2>P_1$.

\begin{enumerate}
\item Représenter le cycle dans le diagramme de Clapeyron $(P,V)$ et préciser la nature de chaque branche.
\item Calculer le travail reçu sur chacune des quatre branches, puis le travail total. Retrouver sa valeur absolue à partir de l'aire du rectangle.
\item Exprimer $\Delta U$ et $Q$ sur chaque branche en fonction de $P_1$, $P_2$, $V_1$ et $V_2$ uniquement.
\item Vérifier que $\Delta U_{\mathrm{cycle}}=0$ et $Q_{\mathrm{cycle}}=-W_{\mathrm{cycle}}$. Le cycle est-il moteur ou récepteur ? Sur quelles branches le gaz reçoit-il de la chaleur ?
\item Application numérique : $P_1=1{,}0$ bar, $P_2=3{,}0$ bar, $V_1=5{,}0$ L et $V_2=15{,}0$ L. Donner $W$ et $Q$ sur chaque branche, puis leurs bilans sur le cycle.
\end{enumerate}

\begin{indication}
Sur une isochore, $nR\,\Delta T=V\,\Delta P$ et le travail est nul. Sur une isobare, $nR\,\Delta T=P\,\Delta V$ et $W=-P\Delta V$. Toute fonction d'état retrouve sa valeur initiale à la fin d'un cycle.
\end{indication}

\begin{solution}
\textbf{Question 1.} Les points $A$, $B$, $C$ et $D$ sont respectivement en bas à gauche, en haut à gauche, en haut à droite et en bas à droite. Les branches $A\to B$ et $C\to D$ sont isochores ; $B\to C$ et $D\to A$ sont isobares. Le parcours est horaire.

\textbf{Questions 2--3.} En posant $\Delta P=P_2-P_1>0$ et $\Delta V=V_2-V_1>0$ :
\[
\begin{array}{c|c|c|c}
\text{Branche} & W & \Delta U & Q=\Delta U-W\\ \hline
A\to B & 0 & \frac32V_1\Delta P & \frac32V_1\Delta P\\[2pt]
B\to C & -P_2\Delta V & \frac32P_2\Delta V & \frac52P_2\Delta V\\[2pt]
C\to D & 0 & -\frac32V_2\Delta P & -\frac32V_2\Delta P\\[2pt]
D\to A & P_1\Delta V & -\frac32P_1\Delta V & -\frac52P_1\Delta V
\end{array}
\]
Le travail total reçu vaut
\[
\boxed{W_{\mathrm{cycle}}=-\Delta P\,\Delta V<0},
\]
dont la valeur absolue est l'aire du rectangle.

\textbf{Question 4.} La somme des variations d'énergie interne est nulle. Le premier principe impose donc
\[
Q_{\mathrm{cycle}}=-W_{\mathrm{cycle}}=\Delta P\,\Delta V>0.
\]
Le gaz fournit globalement du travail : le cycle est moteur. Il reçoit de la chaleur sur $A\to B$ et $B\to C$, puis en cède sur $C\to D$ et $D\to A$.

\textbf{Question 5.} Ici $\Delta P=2{,}0\times10^5$ Pa et $\Delta V=1{,}0\times10^{-2}$ m$^3$ :
\[
\begin{array}{c|r|r}
\text{Branche} & W\ (\mathrm J) & Q\ (\mathrm J)\\ \hline
A\to B & 0 & 1500\\
B\to C & -3000 & 7500\\
C\to D & 0 & -4500\\
D\to A & 1000 & -2500
\end{array}
\]
Ainsi $W_{\mathrm{cycle}}=-2000$ J, $Q_{\mathrm{cycle}}=+2000$ J et $\Delta U_{\mathrm{cycle}}=0$.
\end{solution}
\end{exo}

\begin{exo}[Cycle isobare--isotherme--isochore : un exemple récepteur]{cycle-isobare-isotherme-isochore}
\seoready{true}
\keywords{cycle thermodynamique, isobare, isotherme, isochore, diagramme de Clapeyron, cycle récepteur, gaz diatomique}

Une mole de gaz parfait diatomique ($C_V = \frac{5}{2}R$, $C_P = \frac{7}{2}R$) est initialement dans l'état $A$ de coordonnées $P_0 = 2{,}0\times10^5$ Pa et $V_0 = 14$ L dans le diagramme de Clapeyron. On lui fait subir successivement :
\begin{itemize}
\item une détente isobare $A\to B$ qui triple son volume ;
\item une compression isotherme quasi-statique $B\to C$ qui le ramène au volume initial ;
\item un refroidissement isochore $C\to A$ qui le ramène à l'état initial.
\end{itemize}

\begin{enumerate}
\item Représenter le cycle $ABCA$ dans le diagramme de Clapeyron. Calculer $T_A$, la température $T_B$ de l'isotherme, et la pression maximale $P_C$ atteinte.
\item Calculer $W$ et $Q$ sur chacune des trois branches, puis sur le cycle complet. On remarquera que $RT_A = P_0V_0$.
\item Le cycle ainsi parcouru est-il moteur ou récepteur ? Retrouver ce résultat à partir du sens de parcours dans le diagramme. Comment faudrait-il le parcourir pour en faire un moteur ?
\end{enumerate}

\begin{indication}
$T_A = P_0V_0/R$ et sur l'isobare $T \propto V$. Sur l'isotherme, $\Delta U = 0$ et $W = -nRT\ln(V_f/V_i)$. Le sens antihoraire correspond à un cycle récepteur.
\end{indication}

\begin{solution}
\textbf{Question 1.} $T_A = \dfrac{P_0V_0}{R} = \dfrac{2\times10^5 \times 14\times10^{-3}}{8{,}314} \approx 337$ K. Sur l'isobare, $T_B = 3T_A \approx 1010$ K. Sur l'isotherme $B\to C$ : $P_CV_0 = RT_B = 3P_0V_0$, donc $P_C = 3P_0 = 6{,}0\times10^5$ Pa. Le cycle est bordé par l'isobare $P_0$ en bas, l'isotherme $T_B$ en haut, et l'isochore $V_0$ à gauche.

\textbf{Question 2.} On pose $P_0V_0 = RT_A = 2800$ J.

$A\to B$ (isobare, $\Delta V = 2V_0$) :
\[
W_{AB} = -P_0\cdot 2V_0 = -5600\ \text{J},\quad \Delta U_{AB} = \tfrac{5}{2}R\cdot 2T_A = 5RT_A = 14\,000\ \text{J},\quad Q_{AB} = \tfrac{7}{2}R\cdot2T_A = 19\,600\ \text{J}.
\]

$B\to C$ (isotherme à $3T_A$) : $\Delta U = 0$ et
\[
W_{BC} = -3RT_A\ln\frac{V_0}{3V_0} = 3RT_A\ln 3 \approx +9230\ \text{J}, \qquad Q_{BC} = -W_{BC} \approx -9230\ \text{J}.
\]

$C\to A$ (isochore) : $W_{CA} = 0$ et $Q_{CA} = \Delta U_{CA} = \tfrac{5}{2}R(T_A - 3T_A) = -5RT_A = -14\,000$ J.

Cycle complet : $\Delta U = 0$, $W_{\text{total}} = -5600 + 9230 \approx +3630$ J et $Q_{\text{total}} \approx -3630$ J.

\textbf{Question 3.} $W_{\text{total}} > 0$ : le gaz reçoit du travail et rejette de la chaleur, le cycle est récepteur. C'est cohérent avec le sens de parcours : on parcourt l'isobare (branche basse) de gauche à droite et l'isotherme (branche haute) de droite à gauche, soit un parcours antihoraire. Parcouru dans l'autre sens ($A\to C\to B\to A$), le même cycle fournirait $3630$ J de travail par cycle et deviendrait moteur.
\end{solution}
\end{exo}

\begin{exo}[Compression adiabatique : progressive ou brutale]{compression-progressive-brutale}
\keywords{compression adiabatique, transformation réversible, transformation irréversible, pression extérieure constante, travail, comparaison d'états}

$n$ moles de gaz parfait, de capacités thermiques constantes et de rapport
$\gamma=C_P/C_V>1$, sont enfermées dans un cylindre calorifugé fermé par un
piston mobile sans frottement. L'état initial est
\[
A=(P_0,V_0,T_0).
\]
On veut porter la pression extérieure de $P_0$ à $P_1=xP_0$, avec $x>1$,
puis la ramener à $P_0$.

\begin{enumerate}
\item Faire un schéma du cylindre et représenter la suite des états étudiés :
la compression progressive $A\to B$, puis, dans une autre expérience, la
compression brutale $A\to B'$ suivie du retrait brutal de la surcharge
$B'\to B''$.
\item La pression est augmentée très progressivement de $P_0$ à $P_1$.
Justifier que $A\to B$ est une adiabatique réversible. Déterminer
$P_B$, $T_B$, $V_B$, $W_{AB}$, $Q_{AB}$ et $\Delta U_{AB}$.
\item On repart de $A$ et on impose brutalement la pression extérieure
constante $P_1$. Déterminer complètement l'état final $B'$ ainsi que
$W_{AB'}$, $Q_{AB'}$ et $\Delta U_{AB'}$.
\item Comparer formellement $T_{B'}$ à $T_B$, puis $V_{B'}$ à $V_B$.
Laquelle des deux compressions échauffe le plus le gaz et laquelle le comprime
le plus ?
\item Depuis $B'$, on retire brutalement la surcharge : la pression extérieure
redevient $P_0$. Déterminer l'état final $B''$. Montrer si $T_{B''}$ et
$V_{B''}$ sont plus grands ou plus petits que leurs valeurs initiales.
\item Effectuer le bilan énergétique de $B'\to B''$, puis de l'aller-retour
$A\to B'\to B''$. Expliquer ce que la leçon 6 ajoutera à cette comparaison.
\end{enumerate}

\begin{indication}
Pour la transformation brutale, utiliser
$W=-P_{\mathrm{ext}}(V_f-V_i)$ et seulement l'équation d'état aux états
d'équilibre initial et final. Pour comparer $B$ et $B'$, utiliser l'inégalité
entre moyenne arithmétique pondérée et moyenne géométrique pondérée.
\end{indication}

\begin{solution}
On rappelle que
\[
C_V=\frac{R}{\gamma-1}.
\]

\textbf{Question 1.} Le schéma doit faire apparaître le piston calorifugé et
les deux expériences distinctes :
\[
A\longrightarrow B,
\qquad
A\longrightarrow B'\longrightarrow B''.
\]
Les états $B$ et $B'$ ont la même pression $P_1=xP_0$, tandis que $A$ et
$B''$ ont la même pression $P_0$.

\textbf{Question 2 --- Compression progressive.} L'évolution est
quasi-statique, sans frottement et sans échange thermique : elle est
adiabatique réversible. Les lois de Laplace donnent
\[
P_B=xP_0,
\qquad
T_B=T_0x^{(\gamma-1)/\gamma},
\qquad
V_B=V_0x^{-1/\gamma}.
\]
Comme $Q_{AB}=0$,
\[
W_{AB}=\Delta U_{AB}
=nC_VT_0\left[x^{(\gamma-1)/\gamma}-1\right]>0.
\]

\textbf{Question 3 --- Compression brutale.} Pendant l'évolution, seule la
pression extérieure constante $P_1$ intervient dans le travail :
\[
W_{AB'}=-P_1(V_{B'}-V_0).
\]
À l'équilibre final, $P_{B'}=P_1$. En utilisant
$P_0V_0=nRT_0$, $P_1V_{B'}=nRT_{B'}$, $Q_{AB'}=0$ et
$\Delta U_{AB'}=nC_V(T_{B'}-T_0)$, le premier principe donne
\[
T_{B'}=T_0\frac{1+(\gamma-1)x}{\gamma}.
\]
Posons
\[
\theta=\frac{1+(\gamma-1)x}{\gamma}.
\]
Alors
\[
P_{B'}=xP_0,
\qquad
T_{B'}=\theta T_0,
\qquad
V_{B'}=\frac{\theta}{x}V_0,
\]
et
\[
W_{AB'}=\Delta U_{AB'}
=nC_VT_0(\theta-1)
=\frac{nRT_0}{\gamma}(x-1)>0.
\]

\textbf{Question 4 --- Comparaison.} Le nombre $\theta$ est la moyenne
arithmétique pondérée de $1$ et de $x$. La moyenne géométrique pondérée
associée vaut $x^{(\gamma-1)/\gamma}$. Comme $x>1$,
\[
\theta>x^{(\gamma-1)/\gamma}.
\]
Ainsi,
\[
T_{B'}>T_B.
\]
Les deux états ont la même pression $P_1$ ; l'équation des gaz parfaits impose
donc également
\[
V_{B'}>V_B.
\]
La compression brutale échauffe davantage le gaz, mais le comprime moins que
la compression progressive. Elle lui fournit aussi davantage de travail :
$W_{AB'}>W_{AB}$.

\textbf{Question 5 --- Retrait brutal.} Pour le passage de $P_1$ à $P_0$,
le même calcul donne
\[
\frac{T_{B''}}{T_{B'}}
=\frac{1+(\gamma-1)/x}{\gamma}.
\]
En posant
\[
\phi=\frac{1+(\gamma-1)/x}{\gamma},
\]
on obtient
\[
P_{B''}=P_0,
\qquad
T_{B''}=\theta\phi\,T_0,
\qquad
V_{B''}=\theta\phi\,V_0.
\]
Or
\[
\theta\phi-1
=\frac{\gamma-1}{\gamma^2}
\left(x+\frac1x-2\right)>0
\qquad (x>1).
\]
Par conséquent,
\[
\boxed{T_{B''}>T_0},
\qquad
\boxed{V_{B''}>V_0}.
\]
Après l'aller-retour de la pression extérieure, le gaz ne revient donc pas à
l'état $A$.

\textbf{Question 6 --- Bilans énergétiques.} La détente brutale est
adiabatique :
\[
Q_{B'B''}=0,
\qquad
W_{B'B''}=\Delta U_{B'B''}
=nC_VT_0\theta(\phi-1)<0.
\]
Sur l'aller-retour complet,
\[
Q_{\mathrm{tot}}=0,
\qquad
W_{\mathrm{tot}}=\Delta U_{\mathrm{tot}}
=nC_VT_0(\theta\phi-1)>0.
\]
Les manipulations extérieures ramènent la pression à sa valeur initiale, mais
le gaz a reçu un travail net et reste plus chaud et plus dilaté. La leçon 6
reprendra les mêmes états pour montrer que la compression réversible ne crée
pas d'entropie, tandis que chacune des deux transformations brutales en crée.
\end{solution}
\end{exo}
